\par\medskip
\hypertarget{native-067-article050}{}\pdfbookmark[2]{50.}{article-050}
\begin{center}50.\end{center}
Ratio simplicissima, puncti cuiusvis in superficie sphaerae coelestis situm respectu eclipticae determinandi, fit per ipsius distantiam ab ecliptica (\emph{latitudinem}), atque distantiam puncti, ubi ecliptica a perpendiculo demisso secatur, ab aequinoctio (\emph{longitudinem}). Latitudo, ab utraque eclipticae parte usque ad
\SourceLeaf{068}$90^0$ numerata, in regione boreali ut positiva, in australi ut negativa spectatur. Respondeant corporis coelestis loco heliocentrico, i. e. proiectioni rectae a Sole ad corpus ductae in sphaeram coelestem, longitudo $\lambda$, latitudo $\beta$; sit porro $u$ distantia loci heliocentrici a nodo ascendente (quae \emph{argumentum latitudinis} dicitur), $i$ inclinatio orbitae, $\AscendingNode$ longitudo nodi ascendentis, habebunturque inter $i$, $u$, $\beta$, $\lambda-\AscendingNode$, quae quantitates erunt partes trianguli sphaerici rectanguli, relationes sequentes, quas sine ulla restrictione valere facile evincitur:
\[
\begin{array}{rl}
\mathrm{I.}&\operatorname{tang}(\lambda-\AscendingNode)=\cos i\operatorname{tang}u\\[4pt]
\mathrm{II.}&\operatorname{tang}\beta=\operatorname{tang}i\sin(\lambda-\AscendingNode)\\[4pt]
\mathrm{III.}&\sin\beta=\sin i\sin u\\[4pt]
\mathrm{IV.}&\cos u=\cos\beta\cos(\lambda-\AscendingNode).
\end{array}
\]
Quando $i$ et $u$ sunt quantitates datae, $\lambda-\AscendingNode$ inde per aequ. I determinabitur, ac dein $\beta$ per II vel per III, siquidem $\beta$ non nimis ad $\pm90^0$ appropinquat; formula IV si placet ad calculi confirmationem adhiberi potest. Ceterum formulae I et IV docent, $\lambda-\AscendingNode$ et $u$ semper in eodem quadrante iacere, quoties $i$ est inter $0$ et $90^0$; contra $\lambda-\AscendingNode$ et $360^0-u$ ad eundem quadrantem pertinebunt, quoties $i$ est inter $90^0$ et $180^0$, sive, secundum usum vulgarem, quoties motus est retrogradus; hinc ambiguitas, quam determinatio ipsius $\lambda-\AscendingNode$ per tangentem secundum formulam I relinquit, sponte tollitur.

Formulae sequentes e praecedentium combinatione facile derivantur:
\[
\begin{array}{rl}
\mathrm{V.}&\sin(u-\lambda+\AscendingNode)=2\sin\tfrac{1}{2}i^2\sin u\cos(\lambda-\AscendingNode)\\[4pt]
\mathrm{VI.}&\sin(u-\lambda+\AscendingNode)=\operatorname{tang}\tfrac{1}{2}i\sin\beta\cos(\lambda-\AscendingNode)\\[4pt]
\mathrm{VII.}&\sin(u-\lambda+\AscendingNode)=\operatorname{tang}\tfrac{1}{2}i\operatorname{tang}\beta\cos u\\[4pt]
\mathrm{VIII.}&\sin(u+\lambda-\AscendingNode)=2\cos\tfrac{1}{2}i^2\sin u\cos(\lambda-\AscendingNode)\\[4pt]
\mathrm{IX.}&\sin(u+\lambda-\AscendingNode)=\operatorname{cotang}\tfrac{1}{2}i\sin\beta\cos(\lambda-\AscendingNode)\\[4pt]
\mathrm{X.}&\sin(u+\lambda-\AscendingNode)=\operatorname{cotang}\tfrac{1}{2}i\operatorname{tang}\beta\cos u.
\end{array}
\]
Angulus $u-\lambda+\AscendingNode$, quoties $i$ est infra $90^0$, aut $u+\lambda-\AscendingNode$, quoties $i$ est ultra $90^0$, secundum usum vulgarem \emph{reductio ad eclipticam} dicitur, est scilicet differentia inter longitudinem heliocentricam $\lambda$ atque longitudinem in orbita quae secundum illum usum est $\AscendingNode\pm u$ (secundum nostrum $\AscendingNode+u$). Quoties
\SourceLeaf{069}inclinatio vel parva est vel a $180^0$ parum diversa, ista reductio tamquam quantitas secundi ordinis spectari potest, et in hoc quidem casu praestabit, $\beta$ primo per formulam III ac dein $\lambda$ per VII aut X computare, quo pacto praecisionem maiorem quam per formulam I assequi licebit.

Demisso perpendiculo a loco corporis coelestis in spatio ad planum eclipticae, distantia puncti intersectionis a Sole \emph{distantia curtata} appellatur. Quam per $r'$, radium vectorem autem per $r$ designando, habebimus
\[
\begin{array}{rl}
\mathrm{XI.}&r'=r\cos\beta.
\end{array}
\]
